% =========================================================
% Nearness: Point-Set Topology — Notes Index
% =========================================================

% ---------------------------------------------------------
% Breadcrumb
% ---------------------------------------------------------
\begin{tcolorbox}[
  colback=gray!6,
  colframe=gray!40,
  arc=2pt,
  left=8pt, right=8pt, top=6pt, bottom=6pt,
  title={\small\textbf{Where You Are in the Journey}},
  fonttitle=\small\bfseries
]
\begin{center}
\small
Propositional Logic
$\;\to\;$ Sets \& Functions
$\;\to\;$ $\mathbb{R}$, Real Analysis
$\;\to\;$ Metric Spaces
$\;\to\;$ \textbf{Nearness (Point-Set Topology)}
$\;\to\;$ Finiteness (Compactness)
$\;\to\;$ $\cdots$
\end{center}

\medskip
\noindent\textbf{How we got here.}
Metric spaces showed us that convergence and continuity can be defined
without coordinates --- only distance matters.
Point-set topology takes the final abstraction step: it discards
distance altogether, retaining only the family of open sets
as the primitive notion.

\medskip
\noindent\textbf{What this chapter builds.}
We begin with the topology inherited from metric spaces --- the content
moved here from the metric-spaces chapter --- and then abstract further.
Topics include the open-set framework, nowhere-dense sets, the Baire
Category Theorem, the Generalized Heine-Borel theorem, and the foundations
of compactness and connectedness.

\medskip
\noindent\textbf{Where this leads.}
Compactness is the global finiteness principle that connects topology
back to analysis: continuous functions on compact sets attain their
extrema and are uniformly continuous, unlocking the Riemann integral.
\end{tcolorbox}
\vspace{1em}

% ---------------------------------------------------------
% Structural Role
% ---------------------------------------------------------
\begin{tcolorbox}[
  colback=gray!6,
  colframe=gray!40,
  arc=2pt,
  left=8pt, right=8pt, top=6pt, bottom=6pt,
  title={\small\textbf{Structural Role: Nearness --- Neighborhood Structure Without Distance}},
  fonttitle=\small\bfseries
]
Metric spaces taught us that distance controls convergence.
Nearness --- point-set topology --- goes one step further.

It discards the metric entirely
and retains only the family of open sets:
the neighborhoods that determine local behavior.

Convergence, continuity, compactness, and limit points
can all be defined using only this neighborhood structure.

Topology reveals that many sequence-based results
were secretly statements about open sets.
The abstraction completes a three-stage descent:

\begin{itemize}
  \item Order $\;\to\;$ distance,
  \item Distance $\;\to\;$ neighborhoods,
  \item Neighborhoods $\;\to\;$ open sets.
\end{itemize}

Each stage removes a layer of extra structure
while preserving the phenomena that matter for analysis.
Topology isolates the minimal data needed to discuss
convergence and continuity.
\end{tcolorbox}
\vspace{1em}

% ---------------------------------------------------------
% Structural Roadmap
% ---------------------------------------------------------
\subsection*{Structural Roadmap}

\begin{center}
\textbf{Definitions $\longrightarrow$ Main Theorems
$\longrightarrow$ Consequences and Logical Structure}
\end{center}

The development follows the same architecture used throughout this
project, now abstracting neighborhood structure independently of distance.
The global progression is:

\begin{enumerate}
  \item Open sets, closed sets, and their algebra (from metric spaces)
  \item Nowhere-dense sets and the Baire Category Theorem
  \item Generalized Heine-Borel: compact $\iff$ complete $+$ totally bounded
  \item Compactness via open covers
  \item Sequential compactness and the classical Heine-Borel theorem
  \item Connectedness
\end{enumerate}

\begin{remark*}[Primary sources]
Abbott, \textit{Understanding Analysis}, Chapter~3;
Pugh, \textit{Real Mathematical Analysis}, Chapter~2 \S1--6;
Rudin, \textit{Principles of Mathematical Analysis}, Chapter~2.
\end{remark*}

% ---------------------------------------------------------
% Content
% ---------------------------------------------------------
% =========================================================
% Topology of a Metric Space
% =========================================================

\section{Topology of the Real Metric Space}

% =========================================================
\subsection{Basic Definitions}
% =========================================================


\begin{remark}[Why topology appears here]
The metric space and topological definitions introduced in this section ---
open sets, closed sets, neighborhoods, limit points, and compactness ---
are stated in their natural generality, but their real force only becomes
visible once sequences and limits are in hand.

In particular:
\begin{itemize}
  \item A point $x$ is a \emph{limit point} of a set $A$ if and only if
        some sequence in $A \setminus \{x\}$ converges to $x$.
  \item A set $F$ is \emph{closed} if and only if it contains the limits
        of all convergent sequences in $F$.
  \item A set $K$ is \emph{sequentially compact} if and only if every
        sequence in $K$ has a subsequence converging to a point in $K$.
        By the Heine--Borel theorem, this coincides with compactness for
        subsets of $\mathbb{R}^n$.
\end{itemize}

These equivalences are not incidental --- they reveal that the topological
language and the sequential language are two dialects describing the same
structure. The sequence-based characterizations are generally easier to
work with in proofs, while the open-set definitions generalize more
cleanly to spaces where sequences are insufficient (e.g.\ uncountable
products).

The results that follow --- convergence, Cauchy sequences,
Bolzano--Weierstrass, and subsequence theory --- should be read as
building the sequential side of this correspondence. The topological
interpretations will be noted where they arise.
\end{remark}


\begin{definition}[$\varepsilon$-neighborhood]
Let $x_0\in\mathbb{R}$ and $\varepsilon>0$.
\[
N_\varepsilon(x_0)
:=
\{x\in\mathbb{R} : |x-x_0|<\varepsilon\}.
\]
\end{definition}

\begin{definition}[Open ball]
Let $(X,d)$ be a metric space.
\[
B_\varepsilon(x_0)
:=
\{x\in X : d(x,x_0)<\varepsilon\}.
\]
\end{definition}

\begin{definition}[Closed ball]
\[
\overline{B}_\varepsilon(x_0)
:=
\{x\in X : d(x,x_0)\le\varepsilon\}.
\]
\end{definition}

\begin{definition}[Open set]
A set $U\subseteq X$ is open if
\[
\forall x\in U\;\exists \varepsilon>0
\quad\text{such that}\quad
B_\varepsilon(x)\subseteq U.
\]
\end{definition}

\begin{definition}[Closed set]
A set $F\subseteq X$ is closed if $X\setminus F$ is open.
\end{definition}

\begin{definition}[Open cover]
A family $\{U_\alpha\}_{\alpha\in I}$ of open sets is an open cover of $A\subseteq X$ if
\[
A \subseteq \bigcup_{\alpha\in I} U_\alpha.
\]
\end{definition}

\begin{definition}[Closure]
\[
\overline{A}
=
\bigcap\{F\subseteq X : F \text{ closed and } A\subseteq F\}.
\]
\end{definition}

\begin{definition}[Interior]
\[
A^\circ
=
\bigcup\{U\subseteq X : U \text{ open and } U\subseteq A\}.
\]
\end{definition}

\begin{definition}[Limit point]
\[
\forall \varepsilon>0,\quad
(B_\varepsilon(x)\setminus\{x\})\cap A \neq \varnothing.
\]
\end{definition}

\begin{definition}[Compact set]
Every open cover admits a finite subcover.
\end{definition}

\begin{definition}[Sequential compactness]
Every sequence in $K$ has a convergent subsequence whose limit lies in $K$.
\end{definition}

\begin{definition}[Bounded set]
\[
A \subseteq B_R(x_0)
\quad\text{for some } x_0,R.
\]
\end{definition}

% =========================================================
\subsection{Main Theorems}
% =========================================================

\begin{theorem}[Neighborhood = ball in $\mathbb{R}$]
With $d(x,y)=|x-y|$,
\[
N_\varepsilon(x_0)=B_\varepsilon(x_0).
\]
\end{theorem}

\begin{theorem}[Ball characterization of closure]
\[
x\in \overline{A}
\iff
\forall \varepsilon>0,\;
B_\varepsilon(x)\cap A\neq\varnothing.
\]
\end{theorem}

\begin{theorem}[Ball characterization of interior]
\[
x\in A^\circ
\iff
\exists \varepsilon>0,\;
B_\varepsilon(x)\subseteq A.
\]
\end{theorem}

\begin{theorem}[Heine--Borel Theorem on $\mathbb{R}^n$]
For $K\subseteq\mathbb{R}^n$, the following are equivalent:

\begin{enumerate}
\item $K$ compact
\item $K$ sequentially compact
\item Every infinite subset has a limit point in $K$
\item $K$ closed and bounded
\end{enumerate}

Implication cycle:
\[
(1)\Rightarrow(2)\Rightarrow(3)\Rightarrow(4)\Rightarrow(1).
\]
\end{theorem}



\begin{remark}
Closed and bounded implies compact only in finite-dimensional spaces.
\end{remark}

% =========================================================
\subsection{Geometric Illustration}
% =========================================================

\begin{center}
\begin{tikzpicture}[scale=1.05, line cap=round, line join=round]

\def\rBig{3.0}
\def\rSmall{1.4}

\coordinate (a)  at (0,0);
\coordinate (x0) at (-1.4,0);

\fill[blue!25] (a) circle (\rBig);
\fill[red!30]  (x0) circle (\rSmall);

\draw[blue!70!black, dashed, thick] (a) circle (\rBig);
\draw[red!70!black,  dashed, thick] (x0) circle (\rSmall);

\fill (x0) circle (1.3pt);
\fill (a)  circle (1.3pt);

\node[above left] at (x0) {$x_0$};
\node[above right] at (a) {$a$};

\coordinate (epspt) at ($(x0)+(1.1,0)$);
\draw[thick] (x0) -- (epspt);
\node[above] at ($(x0)!0.55!(epspt)$) {$\varepsilon$};

\node at ($(x0)+(0.2,-0.8)$) {$B(x_0,\varepsilon)$};
\node at ($(a)+(1.6,-1.9)$) {$B(a,r)$};

\end{tikzpicture}
\end{center}

% =========================================================
\subsection{Consequences}
% =========================================================

\begin{remark}[Logical Structure]
\[
\text{Metric}
\Rightarrow
\text{Open Balls}
\Rightarrow
\text{Open Sets}
\Rightarrow
\text{Closure / Interior}
\Rightarrow
\text{Compactness}
\]
\end{remark}

\begin{remark}
In $\mathbb{R}^n$, compactness is equivalent to closed and bounded.
This depends critically on completeness.
\end{remark}

% =========================================================
% Nowhere-Dense Sets and a Preview of Baire's Theorem
% Source: volume-iii/analysis/metric-spaces/notes/metric-spaces/
%
% Dependencies:
%   notes-metric-space.tex       — def:metric
%   notes-open-closed-sets.tex   — def:open, def:closed, def:interior
%   notes-dense-separable.tex    — def:dense
%   notes-completeness.tex       — completeness
% =========================================================

% ---------------------------------------------------------
\section{Nowhere-Dense Sets and Baire's Theorem}
\label{sec:nowhere-dense-baire}
% ---------------------------------------------------------

% ---- Toolkit ----
\begin{tcolorbox}[colback=gray!6, colframe=gray!40, arc=2pt,
  left=6pt, right=6pt, top=4pt, bottom=4pt,
  title={\small\textbf{Nowhere-Dense Sets and Baire --- Quick Reference}},
  fonttitle=\small\bfseries]
\small
\begin{tabular}{l l l l}
\toprule
\textbf{Label} & \textbf{Statement} & \textbf{Proof method} & \textbf{Detail} \\
\midrule
D\ref{def:nowhere-dense}
  & $E$ nowhere-dense iff $\operatorname{int}(\overline{E}) = \varnothing$
  & Definition
  & \hyperref[def:nowhere-dense]{$\downarrow$} \\[4pt]
D\ref{def:meager}
  & Meager $=$ countable union of nowhere-dense sets
  & Definition
  & \hyperref[def:meager]{$\downarrow$} \\[4pt]
P\ref{prop:nowhere-dense-equiv}
  & $E$ nowhere-dense $\iff$ $X \setminus \overline{E}$ is dense
  & Complement characterization
  & \hyperref[prop:nowhere-dense-equiv]{$\downarrow$} \\[4pt]
T\ref{thm:baire-preview}
  & Baire: complete metric space is not meager in itself
  & Statement only (proof deferred)
  & \hyperref[thm:baire-preview]{$\downarrow$} \\[4pt]
\bottomrule
\end{tabular}
\end{tcolorbox}

% ---------------------------------------------------------
\paragraph{Motivation.}
Density gives us a notion of ``large'' subsets: a dense set reaches
every corner of the space.
But we also need a notion of ``small'' or ``thin'' subsets,
one that captures sets which --- no matter how you zoom in ---
never fill any open region.

The integers $\mathbb{Z}$ in $\mathbb{R}$ are a perfect example.
No matter which open interval you look at, $\mathbb{Z}$ does not
fill it; $\mathbb{Z}$ has gaps everywhere.
This is the idea behind \emph{nowhere-dense} sets.

Understanding thin sets leads to one of the deepest structural results
in metric space theory: Baire's theorem, which says that a complete
metric space cannot be built from countably many thin pieces.
We introduce the definitions here and state the theorem;
the proof belongs to the compactness chapter where the key tools are available.

% ---- Definition: Nowhere-Dense ----
\begin{tcolorbox}[colback=defbox, colframe=defborder, arc=2pt,
  left=6pt, right=6pt, top=4pt, bottom=4pt,
  title={\small\textbf{Definition (Nowhere-Dense)}},
  fonttitle=\small\bfseries]
\begin{definition}[Nowhere-Dense]
\label{def:nowhere-dense}
Let $(X,d)$ be a metric space and $E \subseteq X$.
We say $E$ is \emph{nowhere-dense in $X$} if the interior of its closure
is empty:
\[
\operatorname{int}(\overline{E}) = \varnothing.
\]
\end{definition}
\end{tcolorbox}

\begin{remark*}[Fully quantified form]
\[
E \text{ is nowhere-dense}
\;\iff\;
\forall\, x \in \overline{E},\; \forall\, r > 0,\;
B_r(x) \not\subseteq \overline{E}.
\]
No point of $\overline{E}$ has an open ball entirely inside $\overline{E}$.
\end{remark*}

\begin{remark*}[Intuition: nowhere-dense vs.\ not-dense]
These sound similar but are very different.

A set $E$ is \emph{not dense} if there is some open ball it misses entirely:
$\exists\, B_r(x)$ with $B_r(x) \cap E = \varnothing$.

A set $E$ is \emph{nowhere-dense} if \emph{every} open ball contains a
sub-ball that $E$ misses entirely.
Nowhere-dense is a much stronger condition --- it says $E$ fails to be
dense even when restricted to any open region.

Think of nowhere-dense as: ``thin no matter where you zoom in.''
Not-dense just means ``thin somewhere.''
\end{remark*}

\begin{remark*}[Why close before taking interior?]
We use $\operatorname{int}(\overline{E})$ rather than $\operatorname{int}(E)$
because isolated points would otherwise cause confusion.
A single point $\{p\}$ has empty interior but its closure is still $\{p\}$,
which also has empty interior --- so $\{p\}$ is nowhere-dense, as it should be.
The closure step ensures the definition behaves well for any set,
not just closed ones.
For closed sets $E$, the definition simplifies to
$\operatorname{int}(E) = \varnothing$.
\end{remark*}

% ---- Proposition: equivalent characterization ----

\begin{proposition}[Complement characterization of nowhere-dense]
\label{prop:nowhere-dense-equiv}
$E \subseteq X$ is nowhere-dense if and only if $X \setminus \overline{E}$
is dense in $X$.
\end{proposition}

\begin{proof}
By definition, $E$ is nowhere-dense iff $\operatorname{int}(\overline{E}) = \varnothing$.
The interior of $\overline{E}$ is the largest open set contained in
$\overline{E}$, so $\operatorname{int}(\overline{E}) = \varnothing$ iff no
non-empty open set is contained in $\overline{E}$, iff every non-empty open
set meets $X \setminus \overline{E}$.
By Proposition~\ref{prop:dense-equiv}(iii), this is exactly
$\overline{X \setminus \overline{E}} = X$, i.e.\ $X \setminus \overline{E}$
is dense.
\end{proof}

\begin{remark*}[Reading the characterization]
$X \setminus \overline{E}$ dense says: the complement of $\overline{E}$
meets every open set.
In other words, no matter where you look, you can find points
that are not even close to $E$.
This is the precise form of ``$E$ is thin everywhere.''
\end{remark*}

% ---- Examples ----

\begin{example}[Nowhere-dense sets]
\label{ex:nowhere-dense}
\begin{enumerate}[label=(\alph*)]
  \item \textbf{Integers in $\mathbb{R}$.}
        $\mathbb{Z}$ is closed, so $\overline{\mathbb{Z}} = \mathbb{Z}$.
        The interior of $\mathbb{Z}$ is empty (no open interval is contained
        in $\mathbb{Z}$). So $\mathbb{Z}$ is nowhere-dense.

  \item \textbf{The Cantor set $C \subset [0,1]$.}
        $C$ is closed and has empty interior (it contains no intervals).
        So $C$ is nowhere-dense in $\mathbb{R}$,
        despite being uncountable --- a striking demonstration that
        nowhere-dense does not mean small in the sense of cardinality.

  \item \textbf{A single point $\{p\}$.}
        $\overline{\{p\}} = \{p\}$ and $\operatorname{int}(\{p\}) = \varnothing$.
        So every singleton is nowhere-dense.

  \item \textbf{A closed ball's boundary $\partial B_r(x)$ in $\mathbb{R}^n$.}
        The boundary sphere is closed with empty interior, hence nowhere-dense.
\end{enumerate}
\end{example}

\begin{example}[Sets that are NOT nowhere-dense]
\label{ex:not-nowhere-dense}
\begin{enumerate}[label=(\alph*)]
  \item \textbf{$\mathbb{Q}$ in $\mathbb{R}$.}
        $\overline{\mathbb{Q}} = \mathbb{R}$, and
        $\operatorname{int}(\mathbb{R}) = \mathbb{R} \neq \varnothing$.
        So $\mathbb{Q}$ is \emph{not} nowhere-dense,
        even though $\mathbb{Q}$ itself has empty interior.
        (The closure step is essential.)

  \item \textbf{Any open set $U \neq \varnothing$.}
        $\operatorname{int}(\overline{U}) \supseteq \operatorname{int}(U) = U \neq \varnothing$.
        No non-empty open set is nowhere-dense.
\end{enumerate}
\end{example}

% ---- Definition: Meager ----
\begin{tcolorbox}[colback=defbox, colframe=defborder, arc=2pt,
  left=6pt, right=6pt, top=4pt, bottom=4pt,
  title={\small\textbf{Definition (Meager / First Category)}},
  fonttitle=\small\bfseries]
\begin{definition}[Meager Set]
\label{def:meager}
A subset $A \subseteq X$ is \emph{meager} (or \emph{of first category})
if it can be written as a countable union of nowhere-dense sets:
\[
A = \bigcup_{n=1}^{\infty} E_n,
\quad \text{where each } E_n \text{ is nowhere-dense.}
\]
A set that is not meager is called \emph{non-meager} (or
\emph{of second category}).
\end{definition}
\end{tcolorbox}

\begin{remark*}[Intuition: meager as ``topologically small'']
Meager sets are thin in every topological sense.
A countable union of sets each of which is ``thin everywhere'' is
still considered ``thin.''
This gives us a topological analogue of measure-zero sets:
\begin{center}
\begin{tabular}{l l}
\toprule
\textbf{Measure theory} & \textbf{Baire category} \\
\midrule
Measure-zero set & Nowhere-dense set \\
Null set (countable union of measure-zero) & Meager set \\
Sets of positive measure & Non-meager (second category) sets \\
\bottomrule
\end{tabular}
\end{center}
The analogy is imperfect --- a set can be meager and have full measure,
or be non-meager and have measure zero --- but the intuition of
``negligible'' is the right one.
\end{remark*}

\begin{example}[$\mathbb{Q}$ is meager in $\mathbb{R}$]
\label{ex:Q-meager}
Write $\mathbb{Q} = \{q_1, q_2, q_3, \ldots\}$ (countable).
Each singleton $\{q_n\}$ is nowhere-dense.
So $\mathbb{Q} = \bigcup_{n=1}^\infty \{q_n\}$ is meager in $\mathbb{R}$.

This says that the rationals, despite being dense, are
\emph{topologically small}.
Density and meagerness are not opposites: $\mathbb{Q}$ is simultaneously
dense and meager in $\mathbb{R}$.
\end{example}

% ---- Baire's Theorem: Statement only ----

\begin{tcolorbox}[colback=thmbox, colframe=thmborder, arc=2pt,
  left=6pt, right=6pt, top=4pt, bottom=4pt,
  title={\small\textbf{Theorem (Baire Category Theorem --- Abbott \S8.2,
         Pugh Theorem 32)}},
  fonttitle=\small\bfseries]
\begin{theorem}[Baire Category Theorem]
\label{thm:baire-preview}
Let $(X, d)$ be a complete metric space.
Then $X$ is not meager in itself: $X$ cannot be written as a countable
union of nowhere-dense sets.

Equivalently, if $X = \bigcup_{n=1}^\infty E_n$ then at least one $E_n$
has non-empty interior.
\end{theorem}
\end{tcolorbox}

\begin{remark*}[Proof deferred]
The standard proof uses a diagonal/nested-ball argument that requires
completeness in an essential way.
The structure is: choose points from a sequence of shrinking open balls,
each forcing us deeper into an open dense set; completeness ensures the
resulting Cauchy sequence converges to a point inside all of them.
The proof will appear once the compactness tools are in place.
\end{remark*}

\begin{remark*}[What Baire's theorem says about $\mathbb{R}$]
Since $\mathbb{R}$ is complete, it is non-meager: you cannot cover $\mathbb{R}$
by countably many nowhere-dense sets.
But we just showed $\mathbb{Q}$ is meager.
Baire's theorem says the irrationals $\mathbb{R} \setminus \mathbb{Q}$
cannot be meager: if they were, then
$\mathbb{R} = \mathbb{Q} \cup (\mathbb{R} \setminus \mathbb{Q})$
would be the union of two meager sets, hence meager, contradicting
the theorem.

This is a precise sense in which the irrationals are
``much larger'' than the rationals, independent of cardinality:
the irrationals are non-meager, the rationals are meager.
\end{remark*}

\begin{remark*}[Application: most continuous functions are
  nowhere differentiable]
One of the most striking consequences of Baire's theorem is the
following result in $C([0,1])$ with the uniform metric.
Let $D \subseteq C([0,1])$ be the set of functions that are
differentiable at at least one point.
One can show that $D$ is meager in $C([0,1])$.
Since $C([0,1])$ is complete (and hence non-meager by Baire),
$D$ does not cover $C([0,1])$.
Therefore most continuous functions --- in the precise sense of
non-meager complement --- are differentiable \emph{nowhere}.
The differentiable functions, which seem like the ``usual'' ones
from calculus, are in fact the exceptional case.
\end{remark*}

\begin{remark*}[Completeness is essential]
The Baire Category Theorem is false for incomplete spaces.
The rationals $\mathbb{Q}$ (with the usual metric, which is incomplete)
\emph{can} be written as a countable union of singletons, each
nowhere-dense in $\mathbb{Q}$.
So $\mathbb{Q}$ is meager in itself.
Completeness is not just a convenient hypothesis --- it is what prevents
the space from being ``built from nothing.''
\end{remark*}

% =========================================================
% Generalized Heine-Borel Theorem
% Source: volume-v/topology/point-set-topology/notes/
%
% Dependencies:
%   notes-metric-spaces (completeness, totally bounded)
%   notes-topology (compactness via open covers)
%   notes-sequential-characterizations (sequential compactness)
% =========================================================

% ---------------------------------------------------------
\section{Generalized Heine-Borel Theorem}
\label{sec:generalized-heine-borel}
% ---------------------------------------------------------

\begin{tcolorbox}[colback=gray!6, colframe=gray!40, arc=2pt,
  left=6pt, right=6pt, top=4pt, bottom=4pt,
  title={\small\textbf{Generalized Heine-Borel --- Quick Reference}},
  fonttitle=\small\bfseries]
\small
\begin{tabular}{l l l l}
\toprule
\textbf{Label} & \textbf{Statement} & \textbf{Proof method} & \textbf{Detail} \\
\midrule
T\ref{thm:gen-heine-borel}
  & Compact $\iff$ complete $+$ totally bounded
  & Sequential compactness
  & \hyperref[thm:gen-heine-borel]{$\downarrow$} \\[4pt]
C\ref{cor:classical-heine-borel}
  & In $\mathbb{R}^n$: compact $\iff$ closed $+$ bounded
  & Totally bounded $\iff$ bounded in $\mathbb{R}^n$
  & \hyperref[cor:classical-heine-borel]{$\downarrow$} \\[4pt]
\bottomrule
\end{tabular}
\end{tcolorbox}

\paragraph{Motivation.}
The classical Heine-Borel theorem says compact subsets of $\mathbb{R}^n$ are
exactly the closed and bounded ones.
Once you know what compactness and total boundedness mean in general metric
spaces, a natural question arises: which condition is really doing the work?
The answer is total boundedness, not mere boundedness.
The Generalized Heine-Borel theorem isolates the two ingredients ---
completeness and total boundedness --- that together force every sequence
to have a convergent subsequence.

\begin{tcolorbox}[colback=thmbox, colframe=thmborder, arc=2pt,
  left=6pt, right=6pt, top=4pt, bottom=4pt,
  title={\small\textbf{Theorem (Generalized Heine-Borel --- Pugh Theorem 65)}},
  fonttitle=\small\bfseries]
\begin{theorem}[Generalized Heine-Borel]
\label{thm:gen-heine-borel}
Let $(X,d)$ be a complete metric space and $A \subseteq X$.
Then $A$ is compact if and only if $A$ is closed and totally bounded.
\end{theorem}
\end{tcolorbox}

\begin{remark*}[Fully quantified form]
\[
\forall (X,d)\ \text{complete metric space},\;
\forall A \subseteq X:\quad
A\ \text{compact}
\iff
\bigl(A\ \text{closed}\bigr)
\land
\bigl(\forall\,\varepsilon>0,\;
\exists\, n\in\mathbb{N},\;
\exists\, x_1,\ldots,x_n \in X,\;
A \subseteq \bigcup_{k=1}^n B_\varepsilon(x_k)\bigr).
\]
\end{remark*}

\begin{remark*}[Proof strategy]
($\Rightarrow$) Every compact set is closed (as a closed subspace of a complete
space it is complete in its own right) and totally bounded (the open cover by
$\varepsilon$-balls has a finite subcover).

($\Leftarrow$) Let $(a_n)$ be any sequence in $A$.
Since $A$ is totally bounded, Proposition~\ref{prop:totally-bounded-Cauchy}
(in the metric-spaces chapter) produces a Cauchy subsequence $(b_n)$.
Since $X$ is complete, $(b_n)$ converges to some $p \in X$.
Since $A$ is closed, $p \in A$.
So every sequence in $A$ has a subsequence converging in $A$ ---
which is sequential compactness, hence compactness.
The equivalence of sequential compactness and open-cover compactness
in metric spaces is established in the compactness chapter.
\end{remark*}

\begin{tcolorbox}[colback=corbox, colframe=corborder, arc=2pt,
  left=6pt, right=6pt, top=4pt, bottom=4pt,
  title={\small\textbf{Corollary (Classical Heine-Borel in $\mathbb{R}^n$)}},
  fonttitle=\small\bfseries]
\begin{corollary}[Classical Heine-Borel]
\label{cor:classical-heine-borel}
A subset $K \subseteq \mathbb{R}^n$ is compact if and only if
$K$ is closed and bounded.
\end{corollary}
\end{tcolorbox}

\begin{remark*}[Why classical and generalized coincide in $\mathbb{R}^n$]
In $\mathbb{R}^n$, bounded and totally bounded are equivalent
(Proposition~\ref{prop:totally-bounded-Rn} in the metric-spaces chapter).
So in $\mathbb{R}^n$, the Generalized Heine-Borel theorem reads exactly as
the classical one: compact $\iff$ closed $+$ bounded.
The generalized version reveals that it is total boundedness, not bare
boundedness, that is the correct condition in general metric spaces.
\end{remark*}

\begin{remark*}[Why Heine-Borel fails in infinite dimensions]
In $\ell^2$ (the space of square-summable sequences), the closed unit ball
$B = \{x : \|x\|_2 \le 1\}$ is closed and bounded but not compact.
The standard basis vectors $e_n$ satisfy $\|e_m - e_n\|_2 = \sqrt{2}$ for
$m \ne n$, so no subsequence is Cauchy, and $B$ is not sequentially compact.
The unit ball in $\ell^2$ is bounded but not totally bounded, which is the
correct diagnosis: total boundedness, not boundedness, is what fails.

This failure is fundamental to infinite-dimensional analysis and motivates
the study of compact operators in functional analysis.
\end{remark*}

\begin{remark*}[The role of completeness]
The hypothesis that $X$ is complete is essential.
A totally bounded but incomplete space can fail to be compact:
the open interval $(0,1)$ with the usual metric is totally bounded
(it sits inside $[0,1]$ which is compact) but is not complete,
and it is not compact (the sequence $1/n$ has no limit in $(0,1)$).
\end{remark*}


% Future sections (uncomment as content is written):
% \input{volume-v/topology/point-set-topology/notes/notes-compactness}
% \input{volume-v/topology/point-set-topology/notes/notes-connectedness}
